% fig_svd_interlacing.tex
% Singular value bar chart: sigma_k(H), sigma_k(H_tilde), interlacing bounds
% n=6 states (k=1..6); r=2 (dynamic states kept)
% Shows the "spectral gap" between sigma_2 and sigma_3 that controls kappa
\begin{tikzpicture}
\begin{axis}[
  thesis style,
  ybar,
  width  = 0.90\columnwidth,
  height = 5.4cm,
  ylabel={Singular value $\sigma_k$},
  xlabel={Index $k$},
  ymin=0, ymax=105,
  xmin=0.5, xmax=6.5,
  xtick={1,2,3,4,5,6},
  ytick={0,20,40,60,80,100},
  bar width=10pt,
  enlarge x limits=0.08,
  legend pos=north east,
  legend style={font=\scriptsize,cells={anchor=west}},
  grid=both,
  grid style={gray!20,line width=0.3pt},
  clip=false,
]

%% Full model sigma_k(H): 6 singular values, well-separated first two
\addplot[fill=thesisgray!50,draw=thesisgray,bar shift=-7pt]
  coordinates {(1,98.2)(2,61.4)(3,18.7)(4,9.3)(5,4.1)(6,1.8)};
\addlegendentry{$\sigma_k(\mathbf{H})$ (full, $n{=}6$)};

%% Reduced sigma_k(tilde H): only r=2 values; lie within interlacing bounds
\addplot[fill=thesisblue!60,draw=thesisblue,bar shift=7pt]
  coordinates {(1,94.7)(2,58.1)};
\addlegendentry{$\sigma_k(\tilde{\mathbf{H}})$ (reduced, $r{=}2$)};

%% Upper interlacing bound: sigma_k(H)
\addplot[only marks,mark=-,mark size=6pt,thesisred,thick,mark options={line width=1.5pt}]
  coordinates {(1,98.2)(2,61.4)};

%% Lower interlacing bound: sigma_{k+4}(H)
\addplot[only marks,mark=-,mark size=6pt,thesisorange,thick,mark options={line width=1.5pt}]
  coordinates {(1,4.1)(2,1.8)};

%% Error bars connecting upper to lower bounds (vertical lines)
\draw[thesisred,thin,dashed] (axis cs:1,98.2) -- (axis cs:1,4.1);
\draw[thesisred,thin,dashed] (axis cs:2,61.4) -- (axis cs:2,1.8);

%% Spectral gap annotation: sigma_2 vs sigma_3
\draw[thesisblue,<->,thick]
  (axis cs:2.5,18.7) -- node[right,font=\tiny,align=left]{Spectral\\gap $\Delta\sigma$}
  (axis cs:2.5,61.4);
\draw[thesisblue,thin,dotted] (axis cs:2,18.7) -- (axis cs:2.5,18.7);
\draw[thesisblue,thin,dotted] (axis cs:2,61.4) -- (axis cs:2.5,61.4);

%% Condition number annotation for reduced model
\node[thesisblue,font=\tiny,draw=thesisblue!30,fill=white,
      rounded corners=2pt,inner sep=2pt,align=center]
  at (axis cs:1.5,105)
  {$\kappa(\tilde{\mathbf{H}}) = 94.7/58.1 = 1.63 \ll 30$};

%% Reduction boundary
\draw[thesisgray,thin,dashed] (axis cs:2.5,0) -- (axis cs:2.5,105)
  node[above,font=\tiny,thesisgray]{$r{=}2$ cut};

\node[font=\tiny,thesisorange] at (axis cs:1.5,7.5) {lower $\sigma_{k+n_a}$};
\node[font=\tiny,thesisred]    at (axis cs:1.5,100.5) {upper $\sigma_k$};

\end{axis}
\end{tikzpicture}
